%-------------------------
% VitaeTex — fixture de compatibilidad de Jake's Resume
% Adaptación en español y A4 para el spike técnico de Tectonic.
%
% Autor original: Jake Gutierrez
% Basado en: https://github.com/sb2nov/resume
% Fuente: https://github.com/jakegut/resume
% Licencia: MIT; consulte LICENSE-JAKES-RESUME.
%------------------------

\documentclass[a4paper,11pt]{article}

\usepackage{latexsym}
\usepackage[empty]{fullpage}
\usepackage{titlesec}
\usepackage[usenames,dvipsnames]{color}
\usepackage{enumitem}
\usepackage[hidelinks]{hyperref}
\usepackage{fancyhdr}
\usepackage[spanish]{babel}
\usepackage{tabularx}

\pagestyle{fancy}
\fancyhf{}
\fancyfoot{}
\renewcommand{\headrulewidth}{0pt}
\renewcommand{\footrulewidth}{0pt}

\addtolength{\oddsidemargin}{-0.5in}
\addtolength{\evensidemargin}{-0.5in}
\addtolength{\textwidth}{1in}
\addtolength{\topmargin}{-0.5in}
\addtolength{\textheight}{1.0in}

\urlstyle{same}
\raggedbottom
\raggedright
\setlength{\tabcolsep}{0in}

\titleformat{\section}{
  \vspace{-4pt}\scshape\raggedright\large
}{}{0em}{}[\color{black}\titlerule \vspace{-5pt}]

% Tectonic utiliza XeTeX y fuentes OpenType con mapeo Unicode nativo. El
% mecanismo glyphtounicode/pdfgentounicode del original es específico de
% pdfTeX; el script del spike verifica la extracción real del texto.

\newcommand{\resumeItem}[1]{
  \item\small{{#1 \vspace{-2pt}}}
}

\newcommand{\resumeSubheading}[4]{
  \vspace{-2pt}\item
    \begin{tabular*}{0.97\textwidth}[t]{l@{\extracolsep{\fill}}r}
      \textbf{#1} & #2 \\
      \textit{\small #3} & \textit{\small #4} \\
    \end{tabular*}\vspace{-7pt}
}

\newcommand{\resumeProjectHeading}[2]{
  \item
    \begin{tabular*}{0.97\textwidth}{l@{\extracolsep{\fill}}r}
      \small #1 & #2 \\
    \end{tabular*}\vspace{-7pt}
}

\newcommand{\resumeDescription}[1]{
  \item[]\small{{#1 \vspace{-4pt}}}
}

\renewcommand\labelitemii{$\vcenter{\hbox{\tiny$\bullet$}}$}
\newcommand{\resumeSubHeadingListStart}{\begin{itemize}[leftmargin=0.15in,label={}]}
\newcommand{\resumeSubHeadingListEnd}{\end{itemize}}
\newcommand{\resumeItemListStart}{\begin{itemize}}
\newcommand{\resumeItemListEnd}{\end{itemize}\vspace{-5pt}}

\begin{document}

\begin{center}
  \textbf{\Huge \scshape María José Núñez} \\ \vspace{1pt}
  \small Ingeniera de Software $|$ Valparaíso, Chile \\ \vspace{1pt}
  \href{mailto:maria.nunez@example.com}{\underline{maria.nunez@example.com}} $|$
  +56 9 5555 0101 $|$
  \href{https://www.linkedin.com/in/maria-nunez}{\underline{linkedin.com/in/maria-nunez}} $|$
  \href{https://github.com/marianunez}{\underline{github.com/marianunez}}
\end{center}

\section{Perfil profesional}
\small{Ingeniera de software enfocada en productos web mantenibles, accesibles y seguros. Experiencia construyendo sistemas con PHP, TypeScript y PostgreSQL, desde el diseño hasta la operación.}

\section{Educación}
\resumeSubHeadingListStart
  \resumeSubheading
    {Universidad Técnica Federico Santa María}{Valparaíso, Chile}
    {Ingeniería Civil Informática}{Marzo 2014 -- Diciembre 2019}
  \resumeDescription{Formación en estructuras de datos, sistemas distribuidos y seguridad de aplicaciones.}
\resumeSubHeadingListEnd

\section{Experiencia}
\resumeSubHeadingListStart
  \resumeSubheading
    {Ingeniera de Software Senior}{Enero 2023 -- Actualidad}
    {Cooperativa Horizonte}{Santiago, Chile}
  \resumeItemListStart
    \resumeItem{Diseñó servicios web con Laravel y PostgreSQL para procesos internos críticos.}
    \resumeItem{Redujo tiempos de despliegue mediante pruebas automatizadas y contenedores reproducibles.}
  \resumeItemListEnd
  \resumeSubheading
    {Desarrolladora Full Stack}{Marzo 2020 -- Diciembre 2022}
    {Laboratorio Puerto Digital}{Valparaíso, Chile}
  \resumeItemListStart
    \resumeItem{Construyó interfaces accesibles con Vue y TypeScript para equipos distribuidos.}
    \resumeItem{Implementó controles de autorización y trazabilidad sin registrar datos sensibles.}
  \resumeItemListEnd
\resumeSubHeadingListEnd

\section{Proyectos}
\resumeSubHeadingListStart
  \resumeProjectHeading
    {\textbf{VitaeTex} $|$ \emph{Laravel, Vue, TypeScript, PostgreSQL, Tectonic}}
    {Abril 2026 -- Actualidad}
  \resumeItemListStart
    \resumeItem{Aplicación para mantener currículums estructurados y generar documentos LaTeX bajo demanda.}
    \resumeItem{Pipeline aislado que compila documentos sin red y elimina todos los archivos temporales.}
  \resumeItemListEnd
\resumeSubHeadingListEnd

\section{Habilidades técnicas}
\begin{itemize}[leftmargin=0.15in,label={}]
  \small{\item{
    \textbf{Lenguajes}: PHP, TypeScript, SQL \\
    \textbf{Frameworks}: Laravel, Vue, Inertia \\
    \textbf{Datos}: PostgreSQL, modelado relacional \\
    \textbf{Herramientas}: Git, Docker, Tectonic
  }}
\end{itemize}

\section{Certificaciones}
\resumeSubHeadingListStart
  \resumeSubheading
    {Seguridad de aplicaciones web}{Noviembre 2025}
    {Instituto de Tecnología Abierta}{Credencial VITAETEX-2025-001}
\resumeSubHeadingListEnd

\end{document}
